
\documentclass{article}

\usepackage{amsmath}
\usepackage{amssymb}
\usepackage{palatino}
\usepackage{fancyhdr}
\usepackage{fullpage}

% document size parameters

% document parameters
\setlength{\textheight}{9in}

\setlength{\parindent}{0pt}
\setlength{\parskip}{8pt}

\setlength{\parindent}{0.in}
\setlength{\parskip}{0.1in}

\pagestyle{fancyplain}
\setlength{\topmargin}{-0.4in}
\setlength{\headsep}{0.4in}
\lhead{\large{\textbf{MATH 232: Numerical Sol'n of Differential Eq's
II}}--Syllabus}
\rhead{\textit{Spring 2007}}


\begin{document}
\sloppy 
\textbf{Instructor:} 
Prof.~Michael Sprague\\ 
Office: Science \& Engineering 332, 209-228-4179\\
Mobile: (209) 217-7240\\
msprague@ucmerced.edu

\textbf{Lecture Time \& Location:} MW 3:00 pm -- 4:20 pm, COB 209  

\textbf{Discussion Section Time \& Location:} 
F 1:00 pm -- 1:50 pm, COB 209 

\textbf{Office hours:} \\
Monday \& Wednesday:  2:00--2:50 pm  (S\&E 332)\\
Alternatively, students are welcome to e-mail, call, or just stop by my
office.  I am rarely on campus on Tuesdays \& Thursdays, but can usually be
reached via email.

\textbf{Prerequisite:} Math 231 or Consent of Instructor

\textbf{Course Units:} Four

\textbf{Catalog Description:} 
Fundamental methods presented in Math 231 are used as a base for discussing
modern methods for solving partial-differential equations.  Numerical
methods include variational, finite element, collocation,  spectral, and
FFT. Error estimates and implementation issues will be discussed.  A
significant amount of programming will be required.

\textbf{Topics covered:} This is the first semester that this course is
being taught.  As such, it is difficult to predict the amount of material
that will be covered.  However, we will discuss three broad classes of
numerical methods that are very common in science and engineering today:
($i$) finite-difference methods, ($ii$) spectral methods, and ($iii$)
finite-element methods.   It is my hope that students will learn the
fundamental theoretical underpinnings of these methods, and be able to
effectively implement them.   For the latter goal, homework assignments will
have a substantial programming/implementation component.

\textbf{Textbooks:} There are no required textbooks for this course.
However, reading assignments will be taken from the list of titles below,
which are held on a non-circulating basis in the library. If you were to
purchase one book, I would suggest the one by Iserles.
\begin{itemize}

 \item Pinchover \& Rubinstein, \textit{An Introduction to Partial
Differential Equations}, Cambridge University Press, 2005

 \item Iserles, \textit{A First Course in the Numerical Analysis of
Differential Equations}, Cambridge University Press, 2000

\item Canuto et al., \textit{Spectral Methods: Fundamentals in Single
Domains}, Springer, 2006

\end{itemize}

\textbf{Other Useful References:} 
\begin{itemize}

  \item Boyde, \textit{Chebyshev and Fourier Spectral Methods: Second Revised
Edition}, Dover, 2001

  \item  Brenner \& Scott, \textit{The Mathematical Theory of Finite Element
Methods}, Springer, 2002

  \item Gottlieb \& Orszag, \textit{Numerical Analysis of Spectral Methods:
Theory and Applications}, SIAM, 1977

  \item Reddy, \textit{Applied Functional Analysis and Variational Methods
in Engineering}, Krieger, 1991
 
  \item Strikwerda, \textit{Finite Difference Schemes and Partial
Differential Equations}, SIAM, 2004


\end{itemize}

\textbf{Course webpage:} The Math 232 website is part of the UCMCROPS course
management system \mbox{(ucmcrops.ucmerced.edu)}. 
It is available automatically to all students enrolled in
this class. The website contains the course calendar, announcements, and
email list. We will use this site for distributing course materials.

\textbf{Discussion section:}  In general, the discussion section will be
used as a time for group work on assignments.   When we fail to cover
sufficient material in a week's lectures, the Friday discussion section will
be used as a third lecture section.  On several known dates, lecture will be
held in discussion section in place of a missed lecture from that week (see
schedule).

\textbf{Homework, Class Participation \& Grade determination:} The course
grade is based on homework (90\%) and class participation (10\%).  Homework
will be assigned approximately every two weeks.  Late homework will not be
accepted without prior permission.

You are encouraged to work in groups. However, \textbf{all work turned in
must be your own.} You must \textbf{list explicitly any outside sources
employed} (\textit{e.g.}\ websites, \textit{Mathematica}, book other than
the three textbook listed aboe, etc.) for each problem you solve.  This does
not mean that you are allowed to copy a solution should you find it posted
elsewhere (see Academic integrity below).  All programs must be written by
you.

\textbf{Exams:} There will be no exams in this course.

\textbf{Computers \& Software:} You may use a calculator (graphing or
otherwise) or other computational tool (\textit{e.g.}, Mathematica, Maple,
Matlab, etc) to aid in your completion of homework assignments.  

Your homeworks will require a significant amount of programming and data
presentation.   Any language will do, but I recommend Matlab or free
alternative (Octave or python).    Please present data graphically whenever
possible.

\textbf{Dropping the course:}  Please see the UC Merced \textit{General
Catalog} for more details.

\textbf{Special accommodations:} If you qualify for accommodations
because of a disability, please submit a letter from Disability
Services to the instructor in a timely manner so that your needs may
be addressed.  Student Affairs determines accommodations based on
documented disabilities.

I will make every effort to accommodate all students who, because of
religious obligations, have conflicts with scheduled exams, assignments, or
required attendance.  Please speak with me during the first week of classes
regarding any potential academic adjustments or accommodations that may
arise due to religious beliefs during this term.

\clearpage
\textbf{Academic integrity:}  Academic integrity is the foundation of an
academic community and without it none of the educational or research goals
of the university can be achieved.  All members of the university community
are responsible for its academic integrity.  Existing policies forbid
cheating on examinations, plagiarism and other forms of academic dishonesty.
The current policies for UC Merced are described in the  Academic Honesty
Policy (see under Student Judicial Affairs at
http://studentlife.ucmerced.edu) The following general guidelines are
adapted from the UC Davis Code of Academic Conduct
(http://sja.ucdavis.edu/cac.html).

Examples of academic dishonesty include:
\begin{itemize}

\item receiving or providing unauthorized assistance on examinations

\item using unauthorized materials during an examination

\item plagiarism - using materials from sources without citations

\item altering an exam and submitting it for re-grading

\item fabricating data or references

\item using false excuses to obtain extensions of time or to skip
  coursework

\end{itemize}

The ultimate success of a code of academic conduct depends largely on
the degree to which the students fulfill their responsibilities
towards academic integrity.  These responsibilities include:
\begin{itemize}

\item Be honest at all times.

\item Act fairly toward others. For example, do not disrupt or seek an
  unfair advantage over others by cheating, or by talking or allowing
  eyes to wander during exams.

\item Take group as well as individual responsibility for honorable
  behavior.  Collectively, as well as individually, make every effort
  to prevent and avoid academic misconduct, and report acts of
  misconduct which you witness.

\item Do not submit the same work in more than one class. Unless
  otherwise specified by the instructor, all work submitted to fulfill
  course requirements must be work done by the student specifically
  for that course.  This means that work submitted for one course
  cannot be used to satisfy requirements of another course unless the
  student obtains permission from the instructor.

\item Unless permitted by the instructor, do not work with others on
  graded coursework, including in class and take-home tests, papers,
  or homework assignments. When an instructor specifically informs
  students that they may collaborate on work required for a course,
  the extent of the collaboration must not exceed the limits set by
  the instructor.

\item Know what plagiarism is and take steps to avoid it. When using
  the words or ideas of another, even if paraphrased in your own
  words, you must cite your source. Students who are confused about
  whether a particular act constitutes plagiarism should consult the
  instructor who gave the assignment.

\item Know the rules -- ignorance is no defense. 
  Those who violate campus rules regarding academic misconduct are subject to
  disciplinary sanctions, including suspension and dismissal.

\end{itemize}

\newpage


\end{document}
